% !PRINTED none
% Abschnitt-opening recto; no folio is printed.
\begin{center}
\gesperrt{Siebzehnter Abschnitt.}

\textbf{\Large Elliptische Funktionen und quadratische Formen.}

\rule{0.12\linewidth}{0.4pt}
\end{center}

\textbf{§ 114. Singuläre Perioden der doppelt periodischen Funktionen.}

Bezeichnen wir mit $\varphi(u)$ eine doppelt periodische Funktion mit den Perioden $\omega_1$, $\omega_2$, so besitzt, wenn $n$ eine ganze Zahl bedeutet, $\varphi(nu)$ dieselben Perioden, und hierauf beruht die Multiplikation der elliptischen Funktionen, die im sechsten Abschnitt betrachtet wurde. Es entsteht nun die Frage, ob sich nicht noch auf andere Weise ein Multiplikator $\mu$ so bestimmen läßt, daß $\varphi(\mu u)$ die Perioden $\omega_1$, $\omega_2$ besitzt, ob also eine Multiplikation auch mit nicht ganzzahligem Multiplikator existiert. Dies wird dann und nur dann der Fall sein, wenn für ein System ganzer Zahlen $a$, $b$, $c$, $\partial$ die Gleichungen bestehen:
\[
\begin{aligned}
\mu\,\omega_1 &= a\,\omega_1 + b\,\omega_2,\\
\mu\,\omega_2 &= c\,\omega_1 + \partial\,\omega_2.
\end{aligned}
\]
Setzen wir, um die darin enthaltene Forderung näher zu ergründen:
\[
\omega = \dfrac{\omega_2}{\omega_1}
\]
und nehmen an, daß $\omega$ einen positiven imaginären Bestandteil habe, so folgt:
\[\mu=a+b\omega,\]
\weq{1}{\omega=\dfrac{c+\partial\omega}{a+b\omega}.}
Solange $\omega$ als variabel betrachtet wird, ist diese Gleichung nur möglich, wenn $b=c=0$, $a=\partial$ und $\mu$ also eine ganze Zahl ist.
